\section{第一阶段补课：为什么还不能算完成}

前面的第一章已经放进了大量硬件知识，但“知识多”不等于“第一阶段完成”。第一阶段的完成标准不是能背出 CPU、DMA、IOMMU、TLB、PCIe、APIC 这些词，而是能把一个硬件事实推导成一个 OS 设计，再把这个设计落到源码对象、可运行模型、失败案例和观测证据上。若读者看完仍然不知道“为什么必须有这个机制”，第一阶段就没有完成。

\subsection{第一阶段真正要回答的六个问题}

\begin{longtable}{p{0.25\linewidth}p{0.36\linewidth}p{0.27\linewidth}}
\toprule
问题 & 必须能解释到什么程度 & 不合格表现 \\
\midrule
程序为什么会跑 & RIP、取指、译码、寄存器、内存读写、异常 & 只说 CPU 执行指令 \\
OS 为什么能抢回 CPU & 定时器中断、trap frame、内核栈、调度入口 & 只说调度器切换进程 \\
用户程序为什么不能乱改机器 & 特权级、控制寄存器、页权限、MMIO 保护 & 只说内核权限更高 \\
地址为什么会 fault & VMA/PTE/TLB、权限、present、COW、缺页上下文 & 只说内存不够或空指针 \\
设备为什么不是普通变量 & MMIO 副作用、DMA 所有权、中断 completion、屏障 & 只说读写设备寄存器 \\
启动为什么有严格顺序 & 复位向量、固件、DRAM、页表、异常入口、驱动 probe & 只说 bootloader 加载内核 \\
\bottomrule
\end{longtable}

这六个问题是第一阶段的硬标准。后面的进程、文件系统、网络、安全，都要回到这些硬件事实。如果这里没有过关，后面章节会变成一堆互不相连的规则。

\subsection{从单片机到通用 OS 的因果链}

很多教材直接从“进程、线程、虚拟内存”开始，读者就会觉得 OS 是人为规定出来的。更好的方式是从最小单片机一步一步推到通用 OS。

\begin{longtable}{p{0.2\linewidth}p{0.28\linewidth}p{0.38\linewidth}}
\toprule
阶段 & 当前能力 & 下一个机制为什么出现 \\
\midrule
裸循环 & 一个 \code{while(1)} 控制所有设备 & 等待设备会浪费 CPU，需要中断 \\
中断驱动裸机 & 设备能异步通知 CPU & 多个任务仍要合作让出 CPU，需要定时器抢占 \\
抢占式内核 & 定时器可打断死循环 & 任务互不信任，需要特权级和内存保护 \\
多程序系统 & 用户态不能直接控制硬件 & 程序要请求资源，需要系统调用 ABI \\
虚拟内存系统 & 地址空间隔离和懒分配 & 频繁拷贝太贵，需要 COW、mmap、page cache \\
多核系统 & 多任务真正并行 & 共享状态会竞争，需要锁、原子、IPI、per-CPU \\
高速 I/O 系统 & 设备通过 DMA 批量搬运数据 & 生命周期复杂，需要队列、completion、IOMMU \\
\bottomrule
\end{longtable}

注意这个链条里每一步都来自失败场景：轮询浪费 CPU，协作任务可能不让出 CPU，单一地址空间会互相破坏，直接设备访问会越权，PIO 搬数据太慢，多核共享变量会乱序和竞争。OS 不是先验地复杂，而是每个硬件失败场景都迫使它增加一层结构。

\subsection{一条完整路径：键盘字符到用户进程}

用一个小事件把第一阶段串起来：用户按下键盘，终端程序读到一个字符。

\begin{enumerate}
\tightlist
\item 键盘控制器或 USB 主机控制器收到物理事件。
\item 设备把事件放入自己的 FIFO、端点或 DMA buffer。
\item 设备通过中断控制器投递 IRQ 或 MSI。
\item CPU 在指令边界响应中断，硬件保存返回状态，进入内核入口。
\item 内核中断处理读取设备状态，确认事件，必要时把重活交给线程化中断或 workqueue。
\item 驱动把 scancode 或 HID report 转成输入事件，写入内核缓冲区。
\item 等待在 \code{read}、\code{poll} 或 \code{epoll} 上的进程被唤醒。
\item 调度器在某次调度点选择该进程运行。
\item 系统调用返回用户态，用户 buffer 中出现字符或事件结构。
\end{enumerate}

这条路径至少涉及中断、内核栈、设备寄存器、等待队列、调度器、用户指针复制和系统调用返回。若读者只能说“键盘产生中断，程序读到字符”，说明第一阶段仍然没有过关。

\subsection{另一条完整路径：一次文件读到 NVMe}

再看 \code{read(fd, buf, len)} 最终落到 NVMe 的路径。

\begin{lstlisting}
user read(fd, buf, len)
  syscall entry saves user state and checks fd
  VFS resolves file and offset
  page cache lookup misses
  filesystem maps file offset to block
  block layer builds request
  NVMe driver fills submission queue entry
  dma_map user or kernel buffer pages
  driver rings doorbell MMIO register
  device DMA-reads/writes memory
  device posts completion and raises MSI-X
  interrupt handler completes request
  sleeping task wakes and copies data to user buffer
\end{lstlisting}

第一阶段需要读者能指出每个硬件点：doorbell 是 MMIO，数据移动靠 DMA，completion 靠中断或轮询，DMA 地址可能经过 IOMMU，用户 buffer 不能被设备无限期持有，缺页和 page cache 会改变路径。这样进入文件系统章节时，读者看到 page cache、bio、request、completion 才不会断线。

\subsection{为什么必须加可运行模型}

文字解释容易让人以为自己懂了。第一阶段必须至少有一个可运行模型，把下面这些断言跑出来：

\begin{itemize}
\tightlist
\item 用户态不能访问内核页。
\item 用户态不能访问 MMIO。
\item 内核可以通过 MMIO 通知设备。
\item 设备事件会产生中断，CPU 保存 trap frame 并进入内核模式。
\item DMA 映射完成前后有明确所有权，unmap 后设备不能继续写。
\end{itemize}

这些不是实践卷才需要的内容。原理卷也需要最小代码，因为它证明书里的概念能落成状态机。下一节给出一个 C++20 模型，覆盖第一阶段最小硬件契约。
